\section{Análisis}

La \textbf{colección canónica de conjuntos de ítems LR(0)} es la estructura central sobre la
que se construye un analizador sintáctico ascendente determinista. Un \textit{ítem LR(0)} de
la forma $A \to \alpha\,\bullet\,\beta$ representa un estado parcial del reconocimiento:
$\alpha$ ya ha sido desplazado a la pila y se espera ver los símbolos de $\beta$. Cuando el
punto $\bullet$ llega al final ($A \to \alpha\,\bullet$), el ítem indica que la regla
$A \to \alpha$ puede reducirse. El autómata LR(0) utiliza estos ítems para tomar decisiones
de \textit{desplazamiento} (shift) o \textit{reducción} (reduce) sin necesidad de retroceso.

\subsection*{Gramática aumentada}

Antes de construir la colección se \textbf{aumenta} la gramática original: se introduce un
nuevo símbolo inicial $S'$ con la producción única $S' \to S$. Esto garantiza que el estado
de aceptación sea único (el que contiene $S' \to S\,\bullet$) y que el símbolo inicial
original no aparezca a la derecha de ninguna producción, lo que simplifica los algoritmos
de reducción.

\subsection*{Algoritmo de Clausura}

Dado un conjunto de ítems $I$, $\mathrm{CERRADURA}(I)$ se obtiene aplicando iterativamente
la siguiente regla hasta alcanzar el punto fijo:

\begin{enumerate}
  \item Todo ítem de $I$ pertenece a $\mathrm{CERRADURA}(I)$.
  \item Si $A \to \alpha\,\bullet\,B\,\beta \in \mathrm{CERRADURA}(I)$ y
        $B \to \gamma$ es una producción de la gramática,
        entonces $B \to \bullet\,\gamma$ se incorpora a $\mathrm{CERRADURA}(I)$.
\end{enumerate}

La clausura expande el horizonte del autómata: al predecir que puede venir $B$, se agregan
todos los modos posibles de empezar a reconocer $B$. El punto fijo se alcanza cuando no se
puede añadir ningún ítem nuevo.

\subsection*{Algoritmo Ir\_a}

La función $\mathrm{IR\_A}(I, X)$, donde $I$ es un conjunto de ítems y $X$ un símbolo
gramatical, calcula el estado destino al consumir $X$ desde $I$:

\begin{enumerate}
  \item Se forma $J = \{\,A \to \alpha\,X\,\bullet\,\beta \mid
        A \to \alpha\,\bullet\,X\,\beta \in \mathrm{CERRADURA}(I)\,\}$.
  \item El resultado es $\mathrm{CERRADURA}(J)$.
\end{enumerate}

$\mathrm{IR\_A}$ modela la transición del autómata al leer $X$: avanza el punto en todos los
ítems que esperaban $X$ y clausura el resultado para capturar las nuevas predicciones.

\subsection*{Construcción de la colección canónica}

\begin{enumerate}
  \item El estado inicial es $I_0 = \mathrm{CERRADURA}(\{S' \to \bullet\,S\})$.
  \item Para cada estado $I_k$ de la colección y cada símbolo $X$ tal que
        $\mathrm{IR\_A}(I_k, X) \neq \emptyset$:
        \begin{itemize}
          \item Si el estado resultante no existe aún, se añade a la colección y se
                registra la transición $I_k \xrightarrow{X} I_\text{nuevo}$.
          \item Si ya existe, solo se registra la transición.
        \end{itemize}
  \item Se repite hasta que no se generen estados nuevos.
  \item El estado que contiene $S' \to S\,\bullet$ es el estado de \textbf{aceptación}.
\end{enumerate}

\subsection*{Detección de conflictos}

Un autómata LR(0) puede presentar dos tipos de conflictos:

\begin{itemize}
  \item \textbf{Conflicto desplazamiento/reducción}: un estado contiene tanto un ítem con
        el punto antes de un terminal (desplazar) como un ítem con el punto al final
        (reducir).
  \item \textbf{Conflicto reducción/reducción}: un estado contiene dos o más ítems con el
        punto al final, correspondientes a producciones distintas.
\end{itemize}

Una gramática libre de conflictos LR(0) produce tablas de análisis completamente
deterministas. Para gramáticas con conflictos se requieren técnicas más potentes como
SLR(1), LALR(1) o LR(1), que incorporan información de lookahead.

\vspace{0.8cm}

El sistema implementa los tres algoritmos en Python mediante una arquitectura modular
organizada en tres capas: \textit{modelos}, \textit{persistencia} y \textit{presentación}.
La capa de modelos define las estructuras de datos inmutables; la capa de persistencia
agrupa el lector de gramáticas y los algoritmos de clausura, ir\_a y construcción de la
colección; y la capa de presentación contiene la interfaz gráfica construida con Tkinter.
